Long-term conversational memory often requires combining facts recorded
in different sessions and interpreting them relative to a particular
time. Retrieving individually relevant memories does not necessarily
recover the connections needed to answer a question. We introduce
\sys, a framework that formulates memory retrieval as planned witness
search. During memory construction, the system extracts dated
propositions and organizes them into a graph linked to the original
dialogue. During retrieval, an \ac{LLM} uses available evidence to
formulate a conjunctive plan specifying the required relations,
answer form, and temporal reference. An executor searches for facts
that jointly satisfy the plan under consistent variable bindings,
without \ac{LLM} calls. Candidate support can then be checked against
source utterances, while missing relations and rejected candidates
guide bounded replanning. Ranking combines similarity, importance,
and temporal proximity under weights selected for the question.
The reader receives a compact context of dated propositions and
passage summaries, with witnesses guiding evidence selection.
By making relational requirements explicit, \sys connects memory
retrieval to an inspectable account of how stored facts support
an answer.
% Add empirical findings once the evaluation is complete and validated.
